% Generated mechanically from frozen input p01/ko/sets.tex.
% Do not edit this generated file; edit the bound locale source instead.

% OK, start here.
%

\setcounter{chapter}{2}
\chapter{집합론}



\phantomsection
\label{sets-section-phantom}


\section{서론}
\label{sets-section-introduction}

\noindent
때때로 집합론이 필요하다. 여기서는 \cite{Kunen}과 \cite{Jech}에 설명된
선택 공리를 포함한 Zermelo--Fraenkel 집합론(ZFC)을 사용한다.

\section{모든 것은 집합이다}
\label{sets-section-sets-everything}

\noindent
대부분의 수학자는 집합론이 수학의 기본 토대를 제공한다고 생각한다.
그렇다면 이것은 실제로 어떻게 작동하는가? 예를 들어 ``$X$는 스킴이다''라는
문장을 집합론으로 어떻게 옮기는가? 정의를 차례로 풀어 쓰면 된다. 스킴은
각 점이 아핀 스킴인 열린 근방을 갖는 국소환 달린 공간이다. 국소환 달린
공간은 구조층의 모든 줄기가 국소환인 환 달린 공간이다. 환 달린 공간은
위상공간 $X$와 그 위의 환의 층 $\mathcal{O}_X$로 이루어진 쌍
$(X, \mathcal{O}_X)$이다. 위상공간은 집합 $X$와 위상의 공리들을 만족하는
부분집합들의 집합 $\tau \subset \mathcal{P}(X)$로 이루어진 쌍
$(X, \tau)$이다. 이런 식으로 계속한다.

\medskip\noindent
그러면 집합 $S$가 주어졌을 때 그것이 스킴인지 어떻게 판별할 수 있는가?
먼저 집합 $S$가 순서쌍인지 살핀다. 이는 어떤 집합 $a, b$에 대하여 $S$가
$(a, b) := \{\{a\}, \{a, b\}\}$의 꼴이라는 뜻으로 정의된다
(\cite{Jech}, 7쪽 참조). 그렇다면 $a$가 순서쌍 $(c, d)$인지 살핀다.
그렇다면 $d \subset \mathcal{P}(c)$인지 확인하고, 다시 그렇다면 $d$가
집합 $c$ 위의 한 위상을 이루는 집합족인지 확인한다. 이런 식으로 계속한다.

\medskip\noindent
따라서 집합론에서 자유변수 $x$ 하나를 갖고 ``$x$는 스킴이다''라는 개념을
표현하는 완전한 공식 $\phi_{scheme}(x)$을 쓰려면 상당한 작업이 필요하지만,
그렇게 하는 것은 가능하다. 모든 수학적 대상에 대해서도 마찬가지여야 한다.

\section{클래스}
\label{sets-section-classes}

\noindent
비형식적으로는 {\it 클래스}라는 개념을 사용한다. 공식
$\phi(x, p_1, \ldots, p_n)$이 주어졌을 때
$$
C = \{x : \phi(x, p_1, \ldots, p_n)\}
$$
를 {\it 클래스}라고 부른다. 클래스는 그것을 정의하는 공식보다 다루기
쉽지만, 엄밀히 말해 수학적 대상은 아니다. 예를 들어 $R$이 환이면 모든
$R$-가군의 클래스를 생각할 수 있다(결국 ``$M$은 $R$-가군이다''라는 문장을
집합론의 공식으로 옮길 수 있고, 그 공식이 한 클래스를 정의한다).
집합이 아닌 클래스를 {\it 고유 모임}이라고 한다.

\medskip\noindent
이와 같이 $R$-가군의 범주를 생각할 수 있는데, 이는 대상들이 고유 모임을
이루는 ``큰'' 범주이다. 마찬가지로 스킴들의 ``큰'' 범주, 환들의 ``큰'' 범주
등을 생각할 수 있다.



\section{서수}
\label{sets-section-ordinals}

\noindent
$x\in T$이면 $x\subset T$인 집합 $T$를 {\it 추이적}이라고 한다.
집합 $\alpha$가 추이적이고 $\in$에 의해 정렬순서를 이루면 이를
{\it 서수}라고 한다. 이 경우 $\alpha + 1 = \alpha \cup \{\alpha\}$로
정의하며, 이것은 $\alpha$의 {\it 후속}이라 불리는 또 다른 서수이다.
어떤 서수 $\beta$에 대하여 $\alpha = \beta + 1$이면 서수 $\alpha$를
{\it 후속 서수}라고 한다. 가장 작은 서수는 $\emptyset$이며 $0$으로도
표기한다. $\alpha$가 $0$도 아니고 후속 서수도 아니면 $\alpha$를
{\it 극한 서수}라고 하며, 이때
$$
\alpha
=
\bigcup\nolimits_{\gamma \in \alpha} \gamma.
$$
첫 번째 극한 서수는 $\omega$이며, 동시에 첫 번째 무한 서수이다. 첫 번째
비가산 서수 $\omega_1$은 모든 가산 서수의 집합이다. 모든 서수의 모임은
고유 모임이다. 이는 다음 의미에서 $\in$에 의해 정렬되어 있다. 공집합이
아닌 서수의 임의의 집합(더 나아가 클래스)은 최소 원소를 갖는다. 서수들의
집합 $A$가 주어지면 $A$의 {\it 상한}을
$\sup_{\alpha \in A} \alpha =
\bigcup_{\alpha \in A} \alpha$로 정의한다. 이것은 모든
$\alpha \in A$보다 크거나 같은 가장 작은 서수이다. 임의의 정렬집합
$(S, <)$에 대하여 $(S, <) \cong (\alpha, \in)$인 유일한 서수 $\alpha$가
존재하며, 이를 그 정렬집합의 {\it 순서형}이라고 한다.

\section{집합의 계층}
\label{sets-section-sets-hierarchy}

\noindent
초한 재귀로 $V_0 = \emptyset$, $V_{\alpha + 1} = P(V_\alpha)$
(멱집합)로 정의하고, 극한 서수 $\alpha$에 대해서는
$$
V_\alpha = \bigcup\nolimits_{\beta < \alpha} V_\beta.
$$
로 정의한다. 각 $V_\alpha$는 추이적 집합임에 유의하라.

\begin{lemma}
\label{sets-lemma-axiom-regularity}
모든 집합은 어떤 서수 $\alpha$에 대하여 $V_\alpha$의 원소이다.
\end{lemma}

\begin{proof}
\cite[Lemma 6.3]{Jech}을 참조하라.
\end{proof}

\noindent
\cite[Chapter III]{Kunen}에서는 이 보조정리가 기초 공리와 동치임을 설명한다.
집합 $S$의 {\it 랭크}는 $S \in V_{\alpha + 1}$인 가장 작은 서수
$\alpha$이다. {\it 부분 우주}란 문맥에서 크기가 충분히 크다는 것이 분명한
$V_\alpha$ 꼴의 집합을 뜻한다.

\section{기수}
\label{sets-section-cardinals}

\noindent
집합 $A$의 {\it 기수}는 $A$와 전단사인 가장 작은 서수 $\alpha$이다.
이를 나타내기 위해 때때로 $\alpha = |A|$라는 표기를 사용한다. 서수
$\alpha$가 어떤 집합 $A$의 기수로 나타날 때, 다시 말해
$\alpha = |A|$일 때 그리고 그때에만 $\alpha$를 {\it 기수}라고 한다.
기수에는 그리스 문자 $\kappa$, $\lambda$를 사용한다. 첫 번째 무한 기수는
$\omega$이며, 이 문맥에서는 $\aleph_0$로 표기한다. 기수가
$\leq \aleph_0$인 집합을 {\it 가산}이라고 한다. $\alpha$가 서수이면
$\alpha^+$는 $> \alpha$인 가장 작은 기수를 나타낸다. 이를 이용하여
$\aleph_1 = \aleph_0^+$, $\aleph_2 = \aleph_1^+$ 등을 정의할 수 있고,
실제로 임의의 서수 $\alpha$에 대하여 초한 재귀로 $\aleph_\alpha$를
정의할 수 있다. 등식 $\aleph_1 = \omega_1$도 성립한다.

\medskip\noindent
기수 $\kappa, \lambda$의 {\it 덧셈}은 $\kappa \oplus \lambda$로
표기하며, 이는 $\kappa \amalg \lambda$의 기수이다. 기수
$\kappa, \lambda$의 {\it 곱셈}은 $\kappa \otimes \lambda$로 표기하며,
이는 $\kappa \times \lambda$의 기수이다. $\kappa$와 $\lambda$가 무한
기수이면 $\kappa \oplus \lambda = \kappa \otimes \lambda =
\max(\kappa, \lambda)$이다. 기수 $\kappa, \lambda$의 {\it 거듭제곱}은
$\kappa^\lambda$로 표기하며, 이는 $\lambda$에서 $\kappa$로 가는
(집합) 함수들의 집합의 기수이다. 기수들의 임의의 집합 $K$가 주어지면
$K$의 {\it 상한}은
$\sup_{\kappa \in K} \kappa = \bigcup_{\kappa \in K} \kappa$이며,
이것도 기수이다.

\section{공종성}
\label{sets-section-cofinality}

\noindent
정렬집합 $T$의 {\it 공종 부분집합} $S$란
$\forall t \in T \exists s\in S (t \leq s)$를 만족하는 부분집합
$S \subset T$이다. 정렬집합의 부분집합은 유도된 순서에 관하여 정렬집합임에
유의하라. 서수 $\alpha$가 주어졌을 때, $\alpha$의 {\it 공종성}
$\text{cf}(\alpha)$는 $\alpha$의 어떤 공종 부분집합의 순서형으로 나타나는
가장 작은 서수 $\beta$이다. 서수의 공종성은 항상 기수이다. 따라서
$\alpha$의 공종성을 $\alpha$의 공종 부분집합이 가질 수 있는 가장 작은
기수로 정의해도 된다.

\begin{lemma}
\label{sets-lemma-map-from-set-lifts}
$T = \colim_{\alpha < \beta} T_\alpha$가 주어진 서수 $\beta$보다 작은
서수들로 지표화된 집합들의 여극한이라고 하자. 집합의 함수
$\varphi : S \to T$가 주어졌다고 하자. $\beta$가 $S$로 지표화된 서수들의
극한이 아니면, 다시 말해 $\beta$가 $\text{cf}(\beta) > |S|$인 서수이면,
$\varphi$는 어떤 $\alpha < \beta$에 대하여 $T_\alpha$로 가는 함수로
올려진다.
\end{lemma}

\begin{proof}
각 원소 $s \in S$에 대하여 $\alpha_s < \beta$와 $T$에서 $\varphi(s)$로
사상되는 원소 $t_s \in T_{\alpha_s}$를 하나씩 고른다. 가정에 의해
$\alpha = \sup_{s \in S} \alpha_s$는 $\beta$보다 엄밀히 작다. 따라서
$s$에 $T_\alpha$에서의 $t_s$의 상을 대응시키는 함수
$\varphi_\alpha : S \to T_\alpha$가 해가 된다.
\end{proof}

\noindent
다음은 Grothendieck이 특정 아벨 범주에 충분히 많은 단사 대상이 존재함을
증명할 때 사용한, 공종성이 임의로 큰 서수의 존재에 관한 논증과 본질적으로
같다. \cite{Tohoku}를 참조하라.

\begin{proposition}
\label{sets-proposition-exist-ordinals-large-cofinality}
$\kappa$를 기수라 하자. 그러면 공종성이 $\kappa$보다 큰 서수가 존재한다.
\end{proposition}

\begin{proof}
$\kappa$가 유한이면 $\omega = \text{cf}(\omega)$를 택하면 된다.
따라서 $\kappa$가 무한이라고 가정하자. 기수가 $\kappa$보다 엄밀히 큰
가장 작은 서수 $\alpha$를 생각한다. $\text{cf}(\alpha) > \kappa$임을
보이겠다. 만일 $\alpha = \beta + 1$이면 $\alpha$와 $\beta$가 무한이므로
$|\alpha| = |\beta|$이고, 이는 $\alpha$의 최소성에 모순된다. 따라서
$\alpha$는 극한 서수이다. (물론 $\alpha$는 기수이기도 하지만 여기서는
필요하지 않다.) 모순을 얻기 위해 $S \subset \alpha$가
$|S| \leq \kappa$인 공종 부분집합이라고 하자. $\beta \in S$, 즉
$\beta < \alpha$이면 $\alpha$의 최소성에 의해 $|\beta| \leq \kappa$이다.
$\alpha$가 극한 서수이고 $S$가 $\alpha$에서 공종이므로
$\alpha = \bigcup_{\beta \in S} \beta$를 얻는다. 따라서
$|\alpha| \leq |S| \otimes \kappa \leq \kappa \otimes \kappa \leq \kappa$이고,
이는 $\alpha$의 선택에 모순이다.
\end{proof}



\section{반사 원리}
\label{sets-section-reflection-principle}

\noindent
이 내용의 일부는 \cite{Kunen}의 ``쉬운 무모순성 증명''이라는 장에 나온다.

\medskip\noindent
$\phi(x_1, \ldots, x_n)$를 집합론의 공식이라 하자. 이 표기는 $\phi$의 모든
자유변수가 $x_1, \ldots, x_n$ 가운데 나타난다는 것을 뜻한다고 약속하자.
$M$을 집합이라 하자. 공식 $\phi^M(x_1, \ldots, x_n)$은
$\phi(x_1, \ldots, x_n)$에 나오는 모든 $\forall x$와 $\exists x$를 각각
$\forall x\in M$과 $\exists x\in M$으로 바꾸어 얻는 공식이다. 따라서 공식
$\phi(x_1, x_2) = \exists x (x\in x_1 \wedge x\in x_2)$는
$\phi^M(x_1, x_2) = \exists x \in M (x\in x_1 \wedge x\in x_2)$로 바뀐다.
공식 $\phi^M$을 $\phi$의 $M$에 대한 {\it 상대화}라고 한다.

\begin{theorem}
\label{sets-theorem-reflection-principle}
집합론의 공식들로 이루어진 {\bf 유한} 모음
$\phi_1(x_1, \ldots, x_n), \ldots, \phi_m(x_1, \ldots, x_n)$이 주어졌다고
하자. $M_0$를 집합이라 하자. $M_0 \subset M$이고
$\forall x_1, \ldots, x_n \in M$에 대하여
$$
\forall i = 1, \ldots, m, \ \phi_i^{M}(x_1, \ldots, x_n)
\Leftrightarrow
\forall i = 1, \ldots, m, \ \phi_i(x_1, \ldots, x_n).
$$
인 집합 $M$이 존재한다. 실제로 어떤 극한 서수 $\alpha$에 대하여
$M = V_\alpha$로 택할 수 있다.
\end{theorem}

\begin{proof}
\cite[Theorem 12.14]{Jech} 또는 \cite[Theorem 7.4]{Kunen}을 참조하라.
\end{proof}

\noindent
이 정리는 다음과 같이 이해한다. 임의의 $x_1, \ldots, x_n \in M$이 주어졌을
때, 속박변수가 모든 집합을 달리며 공식들이 성립하는 것과 속박변수가
$V_\alpha$의 원소들을 달리며 공식들이 성립하는 것은 동치이다. 이 정리는
집합론의 공식 자체를 직접 다루므로 메타정리이다. 더 구체적으로, 자유변수가
기껏해야 $x_1, \ldots, x_n$인 공식들의 유한 목록
$\phi_1, \ldots, \phi_m$이 주어졌을 때 문장
$$
\begin{matrix}
\forall M_0\ \exists M, \ M_0 \subset M\ \forall x_1, \ldots, x_n \in M \\
\phi_1(x_1, \ldots, x_n) \wedge \ldots \wedge \phi_m(x_1, \ldots, x_n)
\leftrightarrow
\phi_1^M(x_1, \ldots, x_n) \wedge \ldots \wedge \phi_m^M(x_1, \ldots, x_n)
\end{matrix}
$$
이 ZFC에서 증명 가능하다고 말한다. 다시 말해, 실제로 공식들 $\phi_i$의
유한 목록을 적을 때마다 하나의 정리를 얻는다.

\medskip\noindent
공식들 $\phi_i^M(x_1, \ldots, x_n)$의 의미가 그다지 명확하지 않으므로 이
정리를 ``보통 수학''에서 사용하기는 다소 어렵다. 대신 필요한 존재 결과를
직접 증명할 때 반사 원리의 증명에 담긴 생각을 사용할 것이다.

\section{스킴 범주의 구성}
\label{sets-section-categories-schemes}

\noindent
이제 주어진 스킴들의 초기 집합에서 출발하여 자연스러운 연산들의 목록 아래
닫혀 있는 ``작은'' 스킴 범주를 만드는 데 위의 내용을 어떻게 적용하는지
논의한다. 그에 앞서 스킴의 크기를 도입한다. 스킴 $S$가 주어지면
$$
\text{size}(S) = \max(\aleph_0, \kappa_1, \kappa_2),
$$
로 정의하며, 기수 $\kappa_1$과 $\kappa_2$는 다음과 같이 정한다.
\begin{enumerate}
\item $\kappa_1$은 $S$의 아핀 열린집합들의 집합의 기수로 둔다.
\item $\kappa_2$는 모든 아핀 열린집합 $U \subset S$에 대한
모든 $\Gamma(U, \mathcal{O}_S)$의 기수들의 상한으로 둔다.
\end{enumerate}

\begin{lemma}
\label{sets-lemma-bounded-size}
모든 기수 $\kappa$에 대하여 다음을 만족하는 집합 $A$가 존재한다.
$A$의 모든 원소는 스킴이고, $\text{size}(S) \leq \kappa$인 모든 스킴
$S$에 대하여 $X \cong S$인 원소 $X \in A$가 존재한다(스킴의 동형).
\end{lemma}

\begin{proof}
생략한다. 힌트: 임의의 스킴이 아핀 스킴들을 붙여 얻는 스킴과 동형이라는
사실을 어떻게 나타낼지 생각하라.
\end{proof}

\noindent
각 기수 $\kappa$에
\begin{equation}
\label{sets-equation-bound}
Bound(\kappa) = \max\{\kappa^{\aleph_0}, \kappa^+\}.
\end{equation}
를 대응시키는 함수를 $Bound$로 표기한다. 이 함수를 훨씬 빠르게 증가하도록
잡을 수도 있다. 예를 들어 $Bound(\kappa) = \kappa^\kappa$로 두어도 아래의
결과는 여전히 성립한다. 임의의 서수 $\alpha$에 대하여, 대상이
$V_\alpha$의 원소인 스킴 범주의 충만한 부분범주를 $\Sch_\alpha$로 표기한다.
이제 증명할 결과는 다음과 같다.

\begin{lemma}
\label{sets-lemma-construct-category}
위와 같은 표기 $\text{size}$, $Bound$, $\Sch_\alpha$를 사용하자.
$S_0$를 스킴들의 집합이라 하자. 다음 성질들을 갖는 극한 서수 $\alpha$가
존재한다.
\begin{enumerate}
\item
\label{sets-item-inclusion}
$S_0 \subset V_\alpha$이다. 다시 말해
$S_0 \subset \Ob(\Sch_\alpha)$이다.
\item
\label{sets-item-bounded}
임의의 $S \in \Ob(\Sch_\alpha)$와
$\text{size}(T) \leq Bound(\text{size}(S))$인 임의의 스킴 $T$에 대하여,
$T \cong S'$인 스킴 $S' \in \Ob(\Sch_\alpha)$가 존재한다.
\item
\label{sets-item-limit}
임의의 가산\footnote{대상들의 집합과 사상 집합들이 모두 가산이라는
뜻이다. 실제로 (3)과 (4)에서 $\aleph_0$를 임의의 기수로 바꾸어도 보조정리를
증명할 수 있다.} 도표 범주 $\mathcal{I}$와 임의의 함자
$F : \mathcal{I} \to \Sch_\alpha$에 대하여, 극한
$\lim_\mathcal{I} F$가 $\Sch_\alpha$에 존재하는 것과 $\Sch$에 존재하는
것은 동치이며, 이 경우 둘 사이의 자연 사상은 동형이다.
\item
\label{sets-item-colimit}
임의의 가산 지표 범주 $\mathcal{I}$와 임의의 함자
$F : \mathcal{I} \to \Sch_\alpha$에 대하여, 여극한
$\colim_\mathcal{I} F$가 $\Sch_\alpha$에 존재하는 것과 $\Sch$에 존재하는
것은 동치이며, 이 경우 둘 사이의 자연 사상은 동형이다.
\end{enumerate}
\end{lemma}

\begin{proof}
초한 귀납법으로 모든 서수에 서수를 대응시키는 함수 $f$를 다음과 같이
정의한다. $f(0) = 0$으로 둔다. $f(\alpha)$가 주어졌을 때,
$f(\alpha + 1)$을 다음 조건들을 만족하는 가장 작은 서수 $\beta$로 정의한다.
\begin{enumerate}
\item $\alpha + 1 \leq \beta$이고 $f(\alpha) \leq \beta$이다.
\item 임의의 $S \in \Ob(\Sch_{f(\alpha)})$와
$\text{size}(T) \leq Bound(\text{size}(S))$인 임의의 스킴 $T$에 대하여,
$T \cong S'$인 스킴 $S' \in \Ob(\Sch_\beta)$가 존재한다.
\item 임의의 가산 지표 범주 $\mathcal{I}$와 임의의 함자
$F : \mathcal{I} \to \Sch_{f(\alpha)}$에 대하여, 극한
$\lim_\mathcal{I} F$ 또는 여극한 $\colim_\mathcal{I} F$가 $\Sch$에
존재하면 그것은 $\Sch_\beta$의 한 스킴과 동형이다.
\end{enumerate}
그러한 $\beta$가 존재함을 다음과 같이 보인다.
$\Ob(\Sch_{f(\alpha)})$가 집합이므로
\[
\kappa = \sup_{S \in \Ob(\Sch_{f(\alpha)})} Bound(\text{size}(S))
\]
가 존재하며
기수이다. 보조정리 \ref{sets-lemma-bounded-size}에서처럼 $\kappa$에서 출발하여
얻는 스킴들의 집합을 $A$라 하자. 가산 범주들의 집합 $CountCat$이 존재하여,
모든 가산 범주가 $CountCat$의 한 원소와 동형이다. 따라서 위의 (3)에서는
$\mathcal{I}$가 $CountCat$의 원소라고 가정할 수 있다. 이는 (3)의 쌍들
$(\mathcal{I}, F)$이 한 집합 위를 달린다는 뜻이다. 그러므로 (3)과 같은
모든 $(\mathcal{I}, F)$에 대하여 극한이나 여극한이 존재하면 $B$의 한 원소와
동형이 되도록 하는, 스킴들을 원소로 갖는 집합 $B$가 존재한다. 따라서
$A \cup B \subset V_\beta$이고
$\beta > \max\{\alpha + 1, f(\alpha)\}$인 임의의 $\beta$를 택하면
(1)--(3)이 성립한다. 공집합이 아닌 모든 서수 모임은 최소 원소를 가지므로
$f(\alpha + 1)$은 잘 정의된다. 마지막으로 $\alpha$가 극한 서수이면
$f(\alpha) = \sup_{\alpha' < \alpha} f(\alpha')$로 둔다.

\medskip\noindent
$S_0 \subset V_{\beta_0}$인 $\beta_0$를 택하자. 구성에 의해
$f(\beta) \geq \beta$이고, 따라서 $S_0 \subset V_{f(\beta_0)}$이다.
또한 $f$가 비감소이므로 모든 $\beta \geq \beta_0$에 대하여
$S_0 \subset V_{f(\beta)}$이다. 다음으로 $\beta_1 > \beta_0$이고 공종성이
$\text{cf}(\beta_1) > \omega = \aleph_0$인 임의의 서수를 택한다.
서수의 공종성은 임의로 커질 수 있으므로 이는 가능하다. 명제
\ref{sets-proposition-exist-ordinals-large-cofinality}를 참조하라. 보조정리의
문제에 대한 해가 $\alpha = f(\beta_1)$임을 보이겠다.

\medskip\noindent
보조정리의 첫 번째 성질은 위에서 $\beta_1 > \beta_0$로 택했으므로 성립한다.

\medskip\noindent
$\beta_1$의 공종성이 무한이므로 이는 극한 서수이고, 따라서
$f(\beta_1) = \sup_{\beta < \beta_1} f(\beta)$를 얻는다.
그러므로 $\{f(\beta) \mid \beta < \beta_1\} \subset f(\beta_1)$은 공종
부분집합이다. 따라서
$$
V_\alpha = V_{f(\beta_1)} = \bigcup\nolimits_{\beta < \beta_1} V_{f(\beta)}.
$$
이제 $S \in \Ob(\Sch_\alpha)$라 하자. $S \in \Ob(\Sch_{f(\beta)})$인
가장 작은 서수 $\beta$를 $\beta(S)$로 정의한다. 위의 식에 의해 언제나
$\beta(S) < \beta_1$이다. 또한
$\Ob(\Sch_{f(\beta + 1)}) \subset
\Ob(\Sch_\alpha)$이므로, 위의 $f$의 구성에 의해 보조정리의 두 번째 성질이
성립한다.

\medskip\noindent
$\{S_1, S_2, \ldots\} \subset \Ob(\Sch_\alpha)$가 가산 모음이라고 하자.
함수 $\omega \to \beta_1$, $n \mapsto \beta(S_n)$을 생각하자.
$\beta_1$의 공종성이 $> \omega$이므로 이 함수의 상은 공종 부분집합일 수
없다. 따라서 $\{S_1, S_2, \ldots\} \subset \Ob(\Sch_{f(\beta)})$인
$\beta < \beta_1$이 존재한다. 그러므로 임의의 함자
$F : \mathcal{I} \to \Sch_\alpha$는 부분범주들 가운데 하나인
$\Sch_{f(\beta)}$를 통하여 분해된다. 따라서 도표 $F$의 여극한 또는
극한인 스킴 $X$가 존재하면, $f$의 구성에 의해 $X$는 $\Sch_\alpha$의
부분범주인 $\Sch_{f(\beta + 1)}$의 한 대상과 동형이다. 이로써 보조정리의
마지막 두 주장을 증명하였다.
\end{proof}

\begin{remark}
\label{sets-remark-how-to-use-reflection}
위 보조정리는 반사 원리를 사용해서도 증명할 수 있다. 그러나 주의가
필요하다. 문장 $\phi_{scheme}(X)$이 ``$X$는 스킴이다''라는 성질을
표현한다고 할 때, 공식 $\phi_{scheme}^{V_\alpha}(X)$은 무엇을 뜻하는가?
반사 원리에 의해 모든 $X \in V_\alpha$에 대하여
$\phi_{scheme}(X) \leftrightarrow \phi_{scheme}^{V_\alpha}(X)$가 성립하는
$\alpha$를 찾을 수 있는 것은 사실이지만, 이것만으로는 전혀 쓸모가 없다.
그러한 두 문장을 결합해야 비로소 흥미로운 일이 생긴다. 예를 들어
$\phi_{red}(X, Y)$가 ``$X$, $Y$는 스킴이고 $Y$는 $X$의 기약화이다''라는
성질을 표현한다고 하자(스킴, 정의
\ref{schemes-definition-reduced-induced-scheme} 참조). 공식의 쌍
$\phi_1(X, Y) = \phi_{red}(X, Y)$,
$\phi_2(X) = \exists Y, \phi_1(X, Y)$에 반사 원리를 적용한다고 하자.
그러면 반사 원리로 얻은 임의의 $\alpha$는 다음 성질을 갖는다.
$X \in \Ob(\Sch_\alpha)$가 주어지면 $X$의 기약화도
$\Sch_\alpha$의 대상이다(연습문제로 남긴다).
\end{remark}

\begin{lemma}
\label{sets-lemma-bound-affine}
$S$를 아핀 스킴이라 하고 $R = \Gamma(S, \mathcal{O}_S)$라 하자.
그러면 $S$의 크기는 $\max\{ \aleph_0, |R|\}$와 같다.
\end{lemma}

\begin{proof}
$\Spec(R)$의 아핀 열린집합은 기껏해야 $\max\{|R|, \aleph_0\}$개이다.
실제로 임의의 아핀 열린집합 $U \subset \Spec(R)$는 주 열린집합들의 유한 합집합
$D(f_1) \cup \ldots \cup D(f_n)$이므로, 아핀 열린집합의 수는 기껏해야
$\sup_n |R|^n = \max\{|R|, \aleph_0\}$이다. \cite[Ch. I, 10.13]{Kunen}을
참조하라. 한편
$\Gamma(U, \mathcal{O}) \subset R_{f_1} \times \ldots \times R_{f_n}$이므로
$|\Gamma(U, \mathcal{O})| \leq
\max\{\aleph_0, |R_{f_1}|, \ldots, |R_{f_n}|\}$이다. 따라서
$|R_f| \leq \max\{\aleph_0, |R|\}$임을 증명하면 충분하며, 그 증명은 생략한다.
\end{proof}

\begin{lemma}
\label{sets-lemma-bound-size}
$S$를 스킴이라 하고
\[
S = \bigcup_{i \in I} S_i
\]
를 열린 덮개라 하자. 그러면
\[
\text{size}(S) \leq \max\{|I|, \sup_i\{\text{size}(S_i)\}\}
\]
이다.
\end{lemma}

\begin{proof}
$U \subset S$를 임의의 아핀 열린집합이라 하자. $U$는 준콤팩트이므로 유한 개의
원소 $i_1, \ldots, i_n \in I$와 아핀 열린집합
$U_i \subset U \cap S_i$가 존재하여
$U = U_1 \cup U_2 \cup \ldots \cup U_n$이다. 따라서
$$
|\Gamma(U, \mathcal{O}_U)|
\leq
|\Gamma(U_1, \mathcal{O})|
\otimes
\ldots
\otimes
|\Gamma(U_n, \mathcal{O})|
\leq \sup\nolimits_i\{\text{size}(S_i)\}
$$
이다. 또한 이로부터 $S$의 아핀 열린집합들의 집합의 기수가 다음 집합의
기수보다 작거나 같음을 알 수 있다.
$$
\coprod_{n \in \omega}
\coprod_{i_1, \ldots, i_n \in I}
\{\text{affine opens of }S_{i_1}\}
\times
\ldots
\times
\{\text{affine opens of }S_{i_n}\}.
$$
분리합집합 안의 각 집합은 기껏해야 $\sup_i\{\text{size}(S_i)\}$의 기수를
갖는다. 지표집합의 기수는 기껏해야 $\max\{|I|, \aleph_0\}$이다.
\cite[Ch. I, 10.13]{Kunen}을 참조하라. 따라서 \cite[Lemma 5.8]{Jech}에 의해
분리합집합의 기수는 기껏해야 $\max\{\aleph_0, |I|\}
\otimes \sup_i\{\text{size}(S_i)\}$이다. 이로써 보조정리가 성립한다.
\end{proof}

\begin{lemma}
\label{sets-lemma-bound-size-fibre-product}
스킴의 사상
\[
f : X \to S
\]
와
\[
g : Y \to S
\]
가 주어졌다고 하자. 그러면
\[
\text{size}(X \times_S Y) \leq \max\{\text{size}(X), \text{size}(Y)\}
\]
이다.
\end{lemma}

\begin{proof}
$S = \bigcup_{k \in K} S_k$를 아핀 열린 덮개라 하자.
$X = \bigcup_{i \in I} U_i$, $Y = \bigcup_{j \in J} V_j$를 아핀 열린 덮개라
하고, $I$, $J$의 기수는 각각 $\leq \text{size}(X), \text{size}(Y)$라 하자.
각 $i \in I$에 대하여 $f(U_i) \subset \bigcup_{k \in K_i} S_k$인
$k \in K$들의 유한 집합 $K_i$가 존재한다. 각 $j \in J$에 대해서도
$g(V_j) \subset \bigcup_{k \in K_j} S_k$인 $k \in K$들의 유한 집합
$K_j$가 존재한다. 따라서 $f(X), g(Y)$는
$S' = \bigcup_{k \in K'} S_k$에 포함되며, 여기서
$K' = \bigcup_{i \in I} K_i \cup \bigcup_{j \in J} K_j$이다.
$K'$의 기수는 기껏해야 $\max\{\aleph_0, |I|, |J|\}$임에 유의하라.
보조정리 \ref{sets-lemma-bound-size}를 적용하면 모든 $k \in K'$에 대하여
$\text{size}(f^{-1}(S_k) \times_{S_k} g^{-1}(S_k))
\leq \max\{\text{size}(X), \text{size}(Y)\}$임을 보이면 충분하다. 다시 말해
$S$가 아핀이라고 가정해도 된다.

\medskip\noindent
$S$가 아핀이라고 가정하자.
$X = \bigcup_{i \in I} U_i$, $Y = \bigcup_{j \in J} V_j$를 아핀 열린 덮개라
하고, $I$, $J$의 기수는 각각 $\leq \text{size}(X), \text{size}(Y)$라 하자.
다시 보조정리 \ref{sets-lemma-bound-size}에 의해 곱 $U_i \times_S V_j$에 대하여
보조정리를 증명하면 충분하다. 보조정리 \ref{sets-lemma-bound-affine}에 의해
$$
|A \otimes_C B| \leq \max\{\aleph_0, |A|, |B|\}.
$$
임을 보이면 충분하다. 이 부등식의 증명은 생략한다.
\end{proof}

\begin{lemma}
\label{sets-lemma-bound-finite-type}
$S$를 스킴이라 하자. $f : X \to S$가 국소 유한형이고 $X$가 준콤팩트이면
$\text{size}(X) \leq \text{size}(S)$이다.
\end{lemma}

\begin{proof}
각 $U_i$가 $S$의 아핀 열린집합 $S_i$ 안으로 사상되도록 하는 유한 아핀 열린
덮개 $X = \bigcup_{i = 1, \ldots n} U_i$를 찾을 수 있다. 따라서 보조정리
\ref{sets-lemma-bound-size}에 의해 $S$와 $X$가 모두 아핀인 경우로 환원된다.
이 경우 보조정리 \ref{sets-lemma-bound-affine}에 의해
$$
|A[x_1, \ldots, x_n]| \leq \max\{\aleph_0, |A|\}.
$$
임을 보이면 충분하다. 이 부등식의 증명은 생략한다.
\end{proof}

\noindent
대수학, 보조정리 \ref{algebra-lemma-epimorphism-cardinality}에서 환의 전사
$A \to B$에 대하여 $|B| \leq \max(|A|, \aleph_0)$임을 보일 것이다.
스킴에 대한 유사한 결과는 다음 보조정리이다.

\begin{lemma}
\label{sets-lemma-bound-monomorphism}
$f : X \to Y$를 스킴의 단사사상이라 하자. 다음 성질들 중 적어도 하나가
성립하면 $\text{size}(X) \leq \text{size}(Y)$이다.
\begin{enumerate}
\item $f$는 준콤팩트이다.
\item $f$는 국소 유한 표시이다.
% P01-E1042: frozen-source authoring placeholder retained visibly below.
\item 필요에 따라 여기에 더 추가한다.
\end{enumerate}
그러나 국소 유한형인 단사사상에 대해서는 이 상계가 성립하지 않는다.
\end{lemma}

\begin{proof}
$Y = \bigcup_{j \in J} V_j$를 $|J| \leq \text{size}(Y)$인 $Y$의 아핀 열린
덮개라 하자. 보조정리 \ref{sets-lemma-bound-size}에 의해 $X$에서 $V_j$의
역상의 크기를 상계하면 충분하다. 따라서 $Y$가 아핀인 경우, 즉
$Y = \Spec(B)$인 경우로 환원된다. 임의의 아핀 열린집합
$\Spec(A) \subset X$에 대하여
$|A| \leq \max(|B|, \aleph_0) = \text{size}(Y)$이다. 위의 주석과 보조정리
\ref{sets-lemma-bound-affine}를 참조하라. 따라서 $X$의 아핀 열린집합이 기껏해야
$\text{size}(Y)$개임을 보이면 충분하다. $X$가 준콤팩트이면 이는 명백하므로
(1)이 성립한다. (2)의 경우 나타날 수 있는 $B$-대수 $A$의 동형류의 수는
$\text{size}(B)$로 상계된다. 실제로 각 $A$는 $B$ 위의 유한형이므로 어떤
$n, m$ 및 $f_j \in B[x_1, \ldots, x_n]$에 대하여
$B[x_1, \ldots, x_n]/(f_1, \ldots, f_m)$과 동형이다. 그런데
$X \to Y$가 단사사상이므로 $Y = \Spec(B)$ 위의 사상
$\Spec(A) \to X$가 존재한다면 유일하다. 따라서 $X$의 아핀 열린집합의 수는
이들 동형류의 수로 상계된다.

\medskip\noindent
보조정리의 마지막 주장을 증명하기 위해 환
$B = \prod_{n \in \mathbf{N}} \mathbf{F}_2$와 $Y = \Spec(B)$를 생각하자.
$\mathbf{N}$ 위의 각 극대 필터 $\mathcal{U}$에 대하여 잉여체가
$\mathbf{F}_2$인 극대 아이디얼 $\mathfrak m_\mathcal{U}$를 얻는다. 사상
$B \to \mathbf{F}_2$는 원소 $(x_n)$을 $\lim_\mathcal{U} x_n$으로 보낸다.
자세한 내용은 생략한다. 스킴의 사상
$X = \coprod_\mathcal{U} \Spec(\mathbf{F}_2) \to Y$는 모든 점이 서로
다르므로 단사사상이다. 그러나 $X$의 아핀 열린 부분스킴들의 집합의 기수는
$\mathbf{N}$ 위 극대 필터들의 집합의 기수와 같고, 이는
$2^{2^{\aleph_0}}$이다. $|B| = 2^{\aleph_0} < 2^{2^{\aleph_0}}$이므로 결론이
따른다.
\end{proof}

\begin{lemma}
\label{sets-lemma-what-is-in-it}
$\alpha$를 위의 보조정리 \ref{sets-lemma-construct-category}에서와 같은 서수라 하자.
범주 $\Sch_\alpha$는 다음 성질들을 만족한다.
\begin{enumerate}
\item $X, Y, S \in \Ob(\Sch_\alpha)$이면, 임의의 사상
$f : X \to S$, $g : Y \to S$에 대하여 $\Sch_\alpha$ 안의 섬유곱
$X \times_S Y$가 존재하며 스킴 범주에서도 섬유곱이다.
\item $\Ob(\Sch_\alpha)$의 원소들로 이루어진 많아야 가산인 임의의 모음
$S_1, S_2, \ldots$가 주어지면 쌍대곱 $\coprod_i S_i$가
$\Ob(\Sch_\alpha)$에 존재하며 스킴 범주에서도 쌍대곱이다.
\item 임의의 $S \in \Ob(\Sch_\alpha)$와 열린 몰입 $U \to S$에 대하여,
$V \cong U$인 $V \in \Ob(\Sch_\alpha)$가 존재한다.
\item 임의의 $S \in \Ob(\Sch_\alpha)$와 닫힌 몰입 $T \to S$에 대하여,
$S' \cong T$인 $S' \in \Ob(\Sch_\alpha)$가 존재한다.
\item 임의의 $S \in \Ob(\Sch_\alpha)$와 유한형 사상 $T \to S$에 대하여,
$S' \cong T$인 $S' \in \Ob(\Sch_\alpha)$가 존재한다.
\item 스킴 $S$가 다음 조건을 만족하는 열린 덮개
$S = \bigcup_{i \in I} S_i$를 갖는다고 하자. 어떤
$T \in \Ob(\Sch_\alpha)$가 존재하여 (a) 모든 $i \in I$에 대하여
$\text{size}(S_i) \leq \text{size}(T)^{\aleph_0}$이고,
(b) $|I| \leq \text{size}(T)^{\aleph_0}$이다. 그러면 $S$는
$\Sch_\alpha$의 한 대상과 동형이다.
\item 임의의 $S \in \Ob(\Sch_\alpha)$와 국소 유한형 사상
$f : T \to S$에 대하여, $T$가 많아야 $\text{size}(S)^{\aleph_0}$개의
아핀 열린집합으로 덮이면 $S' \cong T$인 $S' \in \Ob(\Sch_\alpha)$가
존재한다. 예를 들어 $T$가 많아야
$|\mathbf{R}| = 2^{\aleph_0} = \aleph_0^{\aleph_0}$개의 아핀 열린집합으로
덮이면 이 조건이 성립한다.
\item 임의의 $S \in \Ob(\Sch_\alpha)$와, 국소 유한 표시이거나
준콤팩트인 단사사상 $T \to S$에 대하여
$S' \cong T$인 $S' \in \Ob(\Sch_\alpha)$가 존재한다.
\item $T \in \Ob(\Sch_\alpha)$가 아핀이라고 하자.
$R = \Gamma(T, \mathcal{O}_T)$로 쓰면, 다음 스킴들은 모두
$\Sch_\alpha$ 안의 한 스킴과 동형이다.
\begin{enumerate}
\item 임의의 아이디얼 $I \subset R$과 완비화
$R^* = \lim_n R/I^n$에 대한 스킴 $\Spec(R^*)$.
\item 임의의 유한형 $R$-대수 $R'$에 대한 스킴 $\Spec(R')$.
\item 임의의 국소화 $S^{-1}R$에 대한 스킴 $\Spec(S^{-1}R)$.
\item 임의의 소 아이디얼 $\mathfrak p \subset R$에 대한 스킴
$\Spec(\overline{\kappa(\mathfrak p)})$.
\item 임의의 부분환 $R' \subset R$에 대한 스킴 $\Spec(R')$.
\item 기수가 기껏해야 $|R|^{\aleph_0}$인 환 위의 임의의 유한형 스킴.
\item 그 밖에도 마찬가지이다.
\end{enumerate}
\end{enumerate}
\end{lemma}

\begin{proof}
(1)과 (2)는 정의에서 바로 따른다. (3)은 $S$의 열린 부분스킴 $U$의 크기가
분명히 $S$의 크기보다 작거나 같으므로 성립한다. (4)는 (5)에서 따르고,
(5)는 (7)에서 따른다. (6)은 보조정리 \ref{sets-lemma-bound-size}에 의해 $S$의
크기가 $\leq \max\{|I|, \sup_i \text{size}(S_i)\} \leq
\text{size}(T)^{\aleph_0}$이므로 성립한다. (7)은 (6)에서 따른다. 실제로
임의의 아핀 열린집합 $V \subset T$에 대하여 보조정리
\ref{sets-lemma-bound-finite-type}에 의해
$\text{size}(V) \leq \text{size}(S)$이다. 따라서 (7)의 상황에 (6)을
적용할 수 있다. (8)은 보조정리 \ref{sets-lemma-bound-monomorphism}에서 따른다.

\medskip\noindent
(9)는 보조정리 \ref{sets-lemma-bound-affine}를 통하여 환
$R^*$, $S^{-1}R$, $\overline{\kappa(\mathfrak p)}$, $R'$ 등의 기수에 대한
상계로 옮겨진다. 아마 가장 흥미로운 것은 환 $R^*$일 것이다. 집합으로서
이는 전사 $R^{\mathbf{N}} \to R^*$의 상이다.
$|R^{\mathbf{N}}| = |R|^{\aleph_0}$이므로, $Bound(\kappa)$를 적어도
$\kappa^{\aleph_0}$로 택한 덕분에 원하는 결론을 얻는다. 휴!
(체의 대수적 폐포의 기수는 그 체의 기수와 같거나 $\aleph_0$이다.)
\end{proof}

\begin{remark}
\label{sets-remark-what-is-not-in-it}
$R$을 환이라 하자. 환 $\prod_{\mathfrak p \in \Spec(R)} \kappa(\mathfrak p)$를
생각한다고 하자. 이 환의 기수는 $|R|^{2^{|R|}}$로 상계되지만 일반적으로
$|R|^{\aleph_0}$로는 상계되지 않는다. 예를 들어
$R = \mathbf{C}[x]$이면 $|R|^{\aleph_0}$로 상계되지 않고,
$R = \prod_{n \in \mathbf{N}} \mathbf{F}_2$이면 $|R|^{|R|}$로 상계되지
않는다. 따라서 위의 보조정리 \ref{sets-lemma-what-is-in-it}에 나오는
``그 밖에도 마찬가지''라는 말은 어느 정도 유보하여 받아들여야 한다.
물론 fppf/\'etale 코호몰로지에 관한 논증에서 이 환들을 고려할 필요가
생긴다면, 위의 함수 $Bound$를 함수 $\kappa \mapsto \kappa^{2^\kappa}$로
바꿀 수 있다.
\end{remark}

\noindent
다음 보조정리에서는 위상론, 절 \ref{topologies-section-fpqc}에서 도입하는
fpqc 덮개의 개념을 사용한다.

\begin{lemma}
\label{sets-lemma-bound-by-covering}
$f : X \to Y$를 스킴의 사상이라 하자. $g_j$가 $f$를 통하여 분해되는
fpqc 덮개 $\{g_j : Y_j \to Y\}_{j \in J}$가 존재한다고 가정하자.
그러면 $\text{size}(Y) \leq \text{size}(X)$이다.
\end{lemma}

\begin{proof}
$V \subset Y$를 아핀 열린집합이라 하자. 정의에 의해 $n \geq 0$과
$a : \{1, \ldots, n\} \to J$, 그리고 아핀 열린집합
$V_i \subset Y_{a(i)}$가 존재하여
$V = g_{a(1)}(V_1) \cup \ldots \cup g_{a(n)}(V_n)$이다.
$f \circ h_j = g_j$인 사상 $h_j : Y_j \to X$를 택하자. 그러면
$h_{a(1)}(V_1) \cup \ldots \cup h_{a(n)}(V_n)$은 $f^{-1}(V)$의 준콤팩트
부분집합이다. 따라서 모든 $i = 1, \ldots, n$에 대하여 $h_{a(i)}(V_i)$를
포함하는 준콤팩트 열린집합 $W \subset f^{-1}(V)$를 찾을 수 있다. 특히
$V = f(W)$이다.

\medskip\noindent
한편 이 사실은 $Y$의 아핀 열린집합들의 집합의 기수가 $X$의 준콤팩트
열린집합들의 집합 $S$의 기수보다 작거나 같음을 보인다. $X$의 모든
준콤팩트 열린집합은 아핀 열린집합들의 유한 합집합이므로 이 집합의 기수는
기껏해야 $\sup |S|^n = \max(\aleph_0, |S|)$이다. 다른 한편
$\{V_i \to V\}$가 fpqc 덮개이므로
$\mathcal{O}_Y(V) \subset \prod_{i = 1, \ldots, n} \mathcal{O}_{Y_{a(i)}}(V_i)$이다.
또한 $V_i \to V$가 $W$를 통하여 분해되므로
$\mathcal{O}_Y(V) \subset \mathcal{O}_X(W)$이다. 다시 $W$는 $X$의 아핀
열린집합들로 이루어진 유한 덮개를 가지므로 $|\mathcal{O}_Y(V)|$는 $X$의
크기로 상계된다. 이제 스킴의 크기 정의에서 보조정리가 따른다.
\end{proof}

\noindent
다음 보조정리에서는 위상론, 절 \ref{topologies-section-fppf}에서 도입하는
fppf 덮개의 개념을 사용한다.

\begin{lemma}
\label{sets-lemma-bound-fppf-covering}
$\{f_i : X_i \to X\}_{i \in I}$를 스킴의 fppf 덮개라 하자.
$\{X_i \to X\}_{i \in I}$의 세분이면서
$\text{size}(\coprod W_j) \leq \text{size}(X)$를 만족하는 fppf 덮개
$\{W_j \to X\}_{j \in J}$가 존재한다.
\end{lemma}

\begin{proof}
$|A| \leq \text{size}(X)$인 아핀 열린 덮개
$X = \bigcup_{a \in A} U_a$를 택하자. 각 $a$에 대하여 유한 부분집합
$I_a \subset I$와, 각 $i \in I_a$에 대하여
$U_a = \bigcup_{i \in I_a} f_i(W_{a, i})$를 만족하는 준콤팩트 열린집합
$W_{a, i} \subset X_i$를 택할 수 있다. 그러면 보조정리
\ref{sets-lemma-bound-finite-type}에 의해
$\text{size}(W_{a, i}) \leq \text{size}(X)$이다. 따라서 보조정리
\ref{sets-lemma-bound-size}에 의해
$\text{size}(\coprod_a \coprod_{i \in I_a} W_{i, a}) \leq \text{size}(X)$이다.
\end{proof}





\section{군 작용을 갖는 집합}
\label{sets-section-sets-with-group-action}

\noindent
$G$를 군이라 하자. $G$-집합들의 ``큰'' 범주를 $G\textit{-Sets}$로 표기한다.
임의의 서수 $\alpha$에 대하여, 대상이 $V_\alpha$에 속하는
$G\textit{-Sets}$의 충만한 부분범주를 $G\textit{-Sets}_\alpha$로 표기한다.
$G$-집합의 크기는
$\text{size}(S) = \max\{\aleph_0, |G|, |S|\}$로 잡는다(여기서 $|G|$와
$|S|$는 바탕 집합들의 기수이다). 위와 같이 함수
$Bound(\kappa) = \kappa^{\aleph_0}$를 사용한다.

\begin{lemma}
\label{sets-lemma-sets-with-group-action}
위와 같은 표기 $G$, $G\textit{-Sets}_\alpha$, $\text{size}$, $Bound$를
사용하자. $S_0$를 $G$-집합들의 집합이라 하자. 다음 성질들을 갖는 극한 서수
$\alpha$가 존재한다.
\begin{enumerate}
\item $S_0 \cup \{{}_GG\} \subset \Ob(G\textit{-Sets}_\alpha)$이다.
\item 임의의 $S \in \Ob(G\textit{-Sets}_\alpha)$와
$\text{size}(T) \leq Bound(\text{size}(S))$인 임의의 $G$-집합 $T$에 대하여,
$T$와 동형인 $S' \in \Ob(G\textit{-Sets}_\alpha)$가 존재한다.
\item 임의의 가산 지표 범주 $\mathcal{I}$와 임의의 함자
$F : \mathcal{I} \to G\textit{-Sets}_\alpha$에 대하여 극한
$\lim_\mathcal{I} F$와 여극한 $\colim_\mathcal{I} F$가
$G\textit{-Sets}_\alpha$에 존재하며, $G\textit{-Sets}$에서의 것과 같다.
\end{enumerate}
\end{lemma}

\begin{proof}
생략한다. 위의 보조정리 \ref{sets-lemma-construct-category}의 증명과 비슷하지만
더 쉽다.
\end{proof}

\begin{lemma}
\label{sets-lemma-what-is-in-it-G-sets}
$\alpha$를 위의 보조정리 \ref{sets-lemma-sets-with-group-action}에서와 같은 서수라
하자. 범주 $G\textit{-Sets}_\alpha$는 다음 성질들을 만족한다.
\begin{enumerate}
\item $G$-집합 ${}_GG$는 $G\textit{-Sets}_\alpha$의 대상이다.
\item (쌍대)곱, 섬유곱, 밂이 $G\textit{-Sets}_\alpha$에 존재하며,
$G\textit{-Sets}$에서의 대응물과 같다.
\item $G\textit{-Sets}_\alpha$의 대상 $U$가 주어졌을 때, 임의의
$G$-안정 부분집합 $O \subset U$는 $G\textit{-Sets}_\alpha$의 한 대상과
동형이다.
\end{enumerate}
\end{lemma}

\begin{proof}
생략한다.
\end{proof}

\section{사이트의 덮개}
\label{sets-section-coverings-site}

\noindent
$\mathcal{C}$가 범주이고(범주론, 정의 \ref{categories-definition-category}),
$\text{Cov}(\mathcal{C})$가 사이트론, 정의 \ref{sites-definition-site}의
성질 (1), (2), (3)을 만족하는 덮개들의 고유 모임이라고 하자. 이 성질들을
다시 적으면 다음과 같다.
\begin{enumerate}
\item $V \to U$가 동형이면
$\{V \to U\} \in \text{Cov}(\mathcal{C})$이다.
\item $\{U_i \to U\}_{i\in I} \in \text{Cov}(\mathcal{C})$이고 각 $i$에
대하여 $\{V_{ij} \to U_i\}_{j\in J_i} \in \text{Cov}(\mathcal{C})$이면
$\{V_{ij} \to U\}_{i \in I, j\in J_i} \in \text{Cov}(\mathcal{C})$이다.
\item $\{U_i \to U\}_{i\in I}\in \text{Cov}(\mathcal{C})$이고
$V \to U$가 $\mathcal{C}$의 사상이면, 모든 $i$에 대하여
$U_i \times_U V$가 존재하며
$\{U_i \times_U V \to V \}_{i\in I} \in \text{Cov}(\mathcal{C})$이다.
\end{enumerate}
서수 $\alpha$에 대하여
$\text{Cov}(\mathcal{C})_\alpha = \text{Cov}(\mathcal{C}) \cap V_\alpha$로
둔다. 서수 $\alpha$와 기수 $\kappa$가 주어지면
$\text{Cov}(\mathcal{C})_{\kappa, \alpha}$를 $|I| \leq \kappa$인 원소
$\mathcal{U} =
\{\varphi_i : U_i \to U\}_{i\in I} \in \text{Cov}(\mathcal{C})_\alpha$들의
집합으로 둔다.

\medskip\noindent
다음 개념을 상기하자. 사이트론, 정의
\ref{sites-definition-combinatorial-tautological}를 참조하라. $\mathcal{C}$에서
같은 공역을 갖는 두 사상족 $\{\varphi_i : U_i \to U\}_{i\in I}$와
$\{\psi_j : W_j \to U\}_{j\in J}$에 대하여, 사상
$\alpha : I \to J$와 $\beta : J\to I$가 존재하여
$\varphi_i = \psi_{\alpha(i)}$ 및 $\psi_j = \varphi_{\beta(j)}$이면 두 사상족을
{\it 조합적으로 동치}라고 한다. 이는 공역이 고정된 사상족들 위의 동치관계를
정의한다.

\begin{lemma}
\label{sets-lemma-coverings-site}
위와 같은 표기를 사용하자.
$\text{Cov}_0 \subset \text{Cov}(\mathcal{C})$를
$\text{Cov}(\mathcal{C})$에 포함되는 집합이라 하자. 다음 성질들을 갖는 기수
$\kappa$와 극한 서수 $\alpha$가 존재한다.
\begin{enumerate}
\item $\text{Cov}_0 \subset \text{Cov}(\mathcal{C})_{\kappa, \alpha}$이다.
\item 덮개들의 집합 $\text{Cov}(\mathcal{C})_{\kappa, \alpha}$는 사이트론,
정의 \ref{sites-definition-site}의 (1), (2), (3)을 만족한다(위를 보라).
다시 말해 $(\mathcal{C}, \text{Cov}(\mathcal{C})_{\kappa, \alpha})$는
사이트이다.
\item $\text{Cov}(\mathcal{C})$의 모든 덮개는
$\text{Cov}(\mathcal{C})_{\kappa, \alpha}$의 한 덮개와 조합적으로 동치이다.
\end{enumerate}
\end{lemma}

\begin{proof}
이를 증명하기 위해 먼저 공역이 고정된 $\mathcal{C}$의 사상 집합들을 모두
모은 집합 $\mathcal{S}$를 생각한다. 다시 말해 $\mathcal{S}$의 원소는
$\text{Arrows}(\mathcal{C})$의 부분집합 $T$로서, $T$의 모든 원소가 같은
공역을 갖는다. 공역이 고정된 사상족
$\mathcal{U} = \{\varphi_i : U_i \to U\}_{i\in I}$가 주어지면
$Supp(\mathcal{U}) = \{ \varphi \in \text{Arrows}(\mathcal{C})
\mid \exists i\in I, \varphi = \varphi_i\}$로 정의한다. 두 사상족
$\mathcal{U} =  \{\varphi_i : U_i \to U\}_{i\in I}$와
$\mathcal{V} = \{V_j \to V\}_{j \in J}$가 조합적으로 동치인 것과
$Supp(\mathcal{U}) = Supp(\mathcal{V})$인 것은 동치임에 유의하라. 다음으로
$\mathcal{S}_\tau \subset \mathcal{S}$를 부분집합
$\mathcal{S}_\tau = \{ T \in \mathcal{S} \mid
\exists\ \mathcal{U} \in \text{Cov}(\mathcal{C}) \ T = Supp(\mathcal{U})\}$로
정의한다. 각 원소 $T \in \mathcal{S}_\tau$에 대하여, 어떤
$\mathcal{U} \in \text{Cov}(\mathcal{C})_\beta$가 존재하여
$T = \text{Supp}(\mathcal{U})$가 되게 하는 가장 작은 서수 $\beta$를
$\beta(T)$로 둔다. 마지막으로
% P01-E0003: frozen undefined S_tau retained in the displayed source formula.
$\beta_0 = \sup_{T \in S_\tau} \beta(T)$로 둔다. 이로부터 모든
$\mathcal{U} \in \text{Cov}(\mathcal{C})$가
$\text{Cov}(\mathcal{C})_{\beta_0}$의 어떤 원소와 조합적으로 동치임이
따른다.

\medskip\noindent
$\kappa$를 $\aleph_0$, 기수 $|\text{Arrows}(\mathcal{C})|$, 그리고
$$
\sup\nolimits_{\{U_i \to U\}_{i\in I} \in \text{Cov}(\mathcal{C})_{\beta_0}}
|I|,
\quad\text{and}\quad
\sup\nolimits_{\{U_i \to U\}_{i\in I} \in \text{Cov}_0} |I|.
$$
의 최댓값으로 둔다. $\kappa$는 무한 기수이므로
$\kappa \otimes \kappa = \kappa$이다. 또한 분명히
$\text{Cov}(\mathcal{C})_{\beta_0} =
\text{Cov}(\mathcal{C})_{\kappa, \beta_0}$이다.

\medskip\noindent
초한 귀납법으로 모든 서수에 서수를 대응시키는 함수 $f$를 다음과 같이
정의한다. $f(0) = 0$으로 둔다. $f(\alpha)$가 주어졌을 때
$f(\alpha + 1)$을 다음 조건들을 만족하는 가장 작은 서수 $\beta$로 정의한다.
\begin{enumerate}
\item $\alpha + 1 \leq \beta$이고 $f(\alpha) \leq \beta$이다.
\item $\{U_i \to U\}_{i\in I}
\in \text{Cov}(\mathcal{C})_{\kappa, f(\alpha)}$이고 각 $i$에 대하여
$\{W_{ij} \to U_i\}_{j\in J_i}
\in \text{Cov}(\mathcal{C})_{\kappa, f(\alpha)}$이면
$\{W_{ij} \to U\}_{i \in I, j\in J_i}
\in \text{Cov}(\mathcal{C})_{\kappa, \beta}$이다.
\item $\{U_i \to U\}_{i\in I}
\in \text{Cov}(\mathcal{C})_{\kappa, \alpha}$이고
$W \to U$가 $\mathcal{C}$의 사상이면
$\{U_i \times_U W \to W \}_{i\in I}
\in \text{Cov}(\mathcal{C})_{\kappa, \beta}$이다.
\end{enumerate}
그러한 $\beta$가 존재함을 보이기 위해, (2)와 (3)에 나타나는 모든 덮개
$\{W_{ij} \to U\}$와 $\{U_i \times_U W \to W \}$의 모임이 분명 집합임에
유의한다. 따라서 이 모든 덮개를 $V_\beta$가 포함하도록 하는 어떤 서수
$\beta$가 존재한다. 또한 덮개 $\{W_{ij} \to U\}$의 지표집합의 기수는
$\sum_{i \in I} |J_i| \leq \kappa \otimes \kappa = \kappa$이므로, 이
덮개들은 $\text{Cov}(\mathcal{C})_{\kappa, \beta}$에 포함된다. 공집합이 아닌
모든 서수 모임은 최소 원소를 가지므로 $f(\alpha + 1)$은 잘 정의된다.
마지막으로 $\alpha$가 극한 서수이면
$f(\alpha) = \sup_{\alpha' < \alpha} f(\alpha')$로 둔다.

\medskip\noindent
다음을 만족하는 서수 $\beta_1$을 택하자.
$\text{Arrows}(\mathcal{C}) \subset V_{\beta_1}$,
% P01-E0004: frozen beta_0 subscript retained below; candidate is sidecar-only.
$\text{Cov}_0 \subset V_{\beta_0}$,
그리고 $\beta_1 \geq \beta_0$이다. 구성에 의해
$f(\beta_1) \geq \beta_1$이고 같은 성질들이 $V_{f(\beta_1)}$에 대해서도
성립한다. 또한 $f$가 비감소이므로 모든 $\beta \geq \beta_1$에 대해서도
그렇다. 다음으로 $\beta_2 > \beta_1$이고 공종성이
$\text{cf}(\beta_2) > \kappa$인 임의의 서수를 택한다. 서수의 공종성은
임의로 커질 수 있으므로 이는 가능하다. 명제
\ref{sets-proposition-exist-ordinals-large-cofinality}를 참조하라. 보조정리의
문제에 대한 해가 $\kappa$, $\alpha = f(\beta_2)$의 쌍임을 보이겠다.

\medskip\noindent
보조정리의 첫 번째와 세 번째 성질은 위에서 $\kappa$와
$\beta_2 > \beta_1 > \beta_0$로 택했으므로 성립한다.

\medskip\noindent
$\beta_2$의 공종성이 무한이므로 이는 극한 서수이고,
$f(\beta_2) = \sup_{\beta < \beta_2} f(\beta)$를 얻는다. 따라서
$\{f(\beta) \mid \beta < \beta_2\} \subset f(\beta_2)$는 공종 부분집합이다.
그러므로
$$
V_\alpha = V_{f(\beta_2)} = \bigcup\nolimits_{\beta < \beta_2} V_{f(\beta)}.
$$
이제 $\mathcal{U} \in \text{Cov}_{\kappa, \alpha}$라 하자.
$\mathcal{U} \in \text{Cov}_{\kappa, f(\beta)}$인 가장 작은 서수 $\beta$를
$\beta(\mathcal{U})$로 정의한다. 위에서 언제나
$\beta(\mathcal{U}) < \beta_2$임을 알 수 있다.

\medskip\noindent
사이트를 정의하는 성질 (1), (2), (3)이 쌍
$(\mathcal{C}, \text{Cov}_{\kappa, \alpha})$에 대하여 성립함을 보여야 한다.
첫 번째 성질은 $\beta_2$의 선택에 의해 $\mathcal{C}$의 모든 화살표가
$V_{f(\beta_2)}$에 포함되므로 성립한다. 세 번째 성질에 대해서는, 덮개
$\mathcal{U} = \{U_i \to U\}_{i \in I}
\in \text{Cov}(\mathcal{C})_{\kappa, \alpha}$가 주어지면
$\beta(\mathcal{U}) < \beta_2$이고 따라서 $\mathcal{U}$의 임의의 밑변환이
$f$의 구성에 의해
% P01-E0005: frozen unbound beta retained in the source formula below.
$\text{Cov}(\mathcal{C})_{\kappa, f(\beta + 1)}$에 포함되며, 따라서
$\text{Cov}(\mathcal{C})_{\kappa, \alpha}$에 포함된다는 사실을 사용한다.

\medskip\noindent
마지막으로 두 번째 조건을 보이자.
% P01-E0006: both frozen f(alpha) stages are retained below.
$\{U_i \to U\}_{i\in I}
\in \text{Cov}(\mathcal{C})_{\kappa, f(\alpha)}$이고 각 $i$에 대하여
$\mathcal{W}_i = \{W_{ij} \to U_i\}_{j\in J_i}
\in \text{Cov}(\mathcal{C})_{\kappa, f(\alpha)}$라고 하자. 함수
$I \to \beta_2$, $i \mapsto \beta(\mathcal{W}_i)$를 생각하자.
$\beta_2$의 공종성은 $> \kappa \geq |I|$이므로 이 함수의 상은 공종
부분집합일 수 없다. 따라서 모든 $i \in I$에 대하여
$\mathcal{W}_i \in \text{Cov}_{\kappa, f(\beta)}$가 되도록 하는
% P01-E0007: frozen beta_1 bound retained below.
$\beta < \beta_1$이 존재한다. 그러므로 덮개
$\{W_{ij} \to U\}_{i\in I, j \in J_i}$는 원하는 대로
$\text{Cov}(\mathcal{C})_{\kappa, f(\beta + 1)}
\subset \text{Cov}(\mathcal{C})_{\kappa, \alpha}$의 원소이다.
\end{proof}

\begin{remark}
\label{sets-remark-better}
어떤 극한 서수 $\alpha$에 대해서는 덮개들의 집합
$\text{Cov}(\mathcal{C})_\alpha$ 자체가 보조정리의 조건들을 만족할 가능성이
크다. 결국 이것이 반사 원리를 적용하면 얻을 법한 결과이다(절
\ref{sets-section-reflection-principle}의 끝과 주석
\ref{sets-remark-how-to-use-reflection}에 설명된 주의사항은 제외하고 말이다).
\end{remark}








\section{아벨 범주와 단사 대상}
\label{sets-section-abelian-categories-injectives}

\noindent
다음 보조정리는 사이트 위 환의 층에 대한 가군들의 범주에 적용된다.

\begin{lemma}
\label{sets-lemma-abelian-injectives}
큰 범주 $\mathcal{A}$가 주어졌다고 하자(범주론, 주석
\ref{categories-remark-big-categories}). $\mathcal{A}$가 아벨 범주이고
충분히 많은 단사 대상을 갖는다고 가정하자. 호몰로지, 정의
\ref{homology-definition-abelian-category}와
\ref{homology-definition-enough-injectives}를 참조하라. 그러면
$\mathcal{A}$의 대상들로 이루어진 임의의 집합 $\{A_s\}_{s\in S}$에 대하여,
다음 성질들을 갖는 아벨 부분범주 $\mathcal{A}' \subset \mathcal{A}$가
존재한다.
\begin{enumerate}
\item 포함 함자 $\mathcal{A}' \to \mathcal{A}$는 완전하다.
\item $\Ob(\mathcal{A}')$은 집합이다.
\item $\Ob(\mathcal{A}')$은 각 $s \in S$에 대한 $A_s$를 포함한다.
\item $\mathcal{A}'$는 충분히 많은 단사 대상을 갖는다.
\item $\mathcal{A}'$의 대상이 단사 대상인 것과 $\mathcal{A}$의 단사 대상인
것은 동치이다.
\end{enumerate}
\end{lemma}

\begin{proof}
생략한다.
\end{proof}
